% =====================================================================
%  APPENDIX A: CALCULATION TRACEABILITY
%  GENERATED by python_scripts/make_appendix_traceability.py (28 Sep 2026)
%  from results/processed/*.json. Do not edit numbers by hand: re-run the script.
% =====================================================================

\chapter{Calculation Traceability}
\label{app:traceability}

This appendix lists every input, formula and result of the analysis in the order the calculations are performed, with the method used for each and the script that performs it \cite{ref32}. The worked values are for the E190; the identical formula set is applied to the A320 and the A350-900 with only the aircraft-specific inputs changed, and their results are given in the tables of Chapters~\ref{ch:concept} to~\ref{ch:results} and in Sections~\ref{sec:app_verification} and~\ref{sec:app_full_matrices}.

Three methods appear in the \textbf{Method} column. \textbf{Direct} is closed-form substitution evaluated once in Python. \textbf{Iterative} is a numerical solution where the equation cannot be rearranged for the unknown. \textbf{FEM} marks quantities that no formula produces: Abaqus solves the discretised equilibrium and heat balance and the result is read from the output database.

\section{Input Values}
\label{sec:app_inputs}

\begin{table}[!htbp]
\centering
\caption{Input values and their sources}
\label{tab:app_inputs}
\footnotesize
\setlength{\tabcolsep}{4pt}
\begin{tabularx}{\textwidth}{@{}llrl>{\raggedright\arraybackslash}X@{}}
\toprule
\textbf{Symbol} & \textbf{Quantity} & \textbf{Value} & \textbf{Unit} & \textbf{Source} \\
\midrule
\multicolumn{5}{@{}l}{\textit{Material: aluminium 2219-T87}} \\
$E$ & Elastic modulus & 73.1 & GPa & MIL-HDBK-5J \cite{ref30} \\
$\nu$ & Poisson's ratio & 0.33 & - & MIL-HDBK-5J \cite{ref30} \\
$\alpha$ & Thermal expansion coefficient & $22.3\times10^{-6}$ & 1/K & MIL-HDBK-5J \cite{ref30} \\
$\rho$ & Density & 2850 & kg/m$^3$ & MIL-HDBK-5J \cite{ref30} \\
$\kappa$ & Thermal conductivity & 116 & W/m$\cdot$K & MIL-HDBK-5J \cite{ref30} \\
$F_{ty}$ & Yield strength, room temperature & 393 & MPa & MIL-HDBK-5J \cite{ref30}; room temperature used deliberately, Section~\ref{sec:concept_material} \\
\addlinespace
\multicolumn{5}{@{}l}{\textit{Fluid and operating condition}} \\
$\rho_{LH2}$ & Liquid hydrogen density & 70.8 & kg/m$^3$ & Leachman et al.\ \cite{ref14} \\
$p$ & Strength design differential & 1.5 & bar & Cruise differential with margin, Section~\ref{ssec:boiloff_requirement} \\
$p_{\min}$ & Stability design differential & 0.187 & bar & $1.2 - 1.01325$ bar, Section~\ref{sec:failure_criteria} \\
$T_{liq}$ & Liquid and ullage sink temperature & 20.3 & K & Bagarello et al.\ \cite{ref29} \\
$T_{ref}$ & Stress-free reference temperature & 20.3 & K & Cold operating state, Section~\ref{sec:coupling} \\
$h_{wet}$, $h_{dry}$ & Film coefficients & 1000, 1.46 & W/m$^2\cdot$K & Bagarello et al.\ \cite{ref29} \\
$LHV_{JetA}$, $LHV_{LH2}$ & Lower heating values & 43, 120 & MJ/kg & Tiwari et al.\ \cite{ref2} \\
\addlinespace
\multicolumn{5}{@{}l}{\textit{Rules fixed across all three aircraft}} \\
$2R/D$ & Tank-to-fuselage diameter ratio & 0.8648 & - & Clearance rule, Section~\ref{sec:concept_geometry} \\
$a/b$ & Cap aspect ratio & 1.6 & - & Section~\ref{sec:concept_geometry}; as Bagarello et al.\ \cite{ref29} \\
$f_{worst}$ & Structural design fill & 0.9 & - & Sizes the volume, Section~\ref{ssec:boiloff_requirement} \\
$f_{max}$ & Operational maximum fill & 0.97 & - & Heat leak budget only, Section~\ref{ssec:boiloff_requirement} \\
SF & Safety factor on yield & 1.5 & - & CS-25.303 \cite{ref22} \\
FOS$_b$ & Target factor on buckling & 1.5 & - & Section~\ref{sec:failure_criteria} \\
$t_{hold}$ & No-vent hold time & 14.7 & h & Adopted budget, Section~\ref{ssec:boiloff_requirement} \\
$k_{eff}$ & MLI effective conductivity & 0.0003 & W/m$\cdot$K & Adopted engineering value, Section~\ref{ssec:insulation_selection} \\
\addlinespace
\multicolumn{5}{@{}l}{\textit{Aircraft-specific, E190}} \\
$D$ & Fuselage outer diameter & 3.01 & m & Reference aircraft \\
$m_{JetA}$ & Kerosene load replaced & 14\,243 & kg & Design requirement, Table~\ref{tab:jeta_basis} \\
\bottomrule
\end{tabularx}
\end{table}


\section{Stage 1: Geometry from the Mission}
\label{sec:app_geometry}

\begin{table}[!htbp]
\centering
\caption{Geometry sizing calculations, E190. Script: final\_numbers.py.}
\label{tab:app_geometry}
\footnotesize
\setlength{\tabcolsep}{4pt}
\begin{tabularx}{\textwidth}{@{}c>{\raggedright\arraybackslash}p{2.6cm}>{\raggedright\arraybackslash}X>{\raggedright\arraybackslash}p{2.4cm}rl@{}}
\toprule
\textbf{\#} & \textbf{Computes} & \textbf{Formula} & \textbf{Inputs} & \textbf{Result} & \textbf{Method} \\
\midrule
1 & Hydrogen mass for the same mission energy & $m_{LH2} = m_{JetA}\,LHV_{JetA}/LHV_{LH2}$ & 14\,243, 43, 120 & 5103.8 kg & Direct \\
\addlinespace
2 & Tank radius & $R = 0.8648\,D/2$ & 3.01 m & 1.3015 m & Direct \\
\addlinespace
3 & Cap height & $h_{cap} = R/1.6$ & $R$ & 0.8134 m & Direct \\
\addlinespace
4 & Required volume & $V = m_{LH2}/(\rho_{LH2}\,f_{worst})$ & 5103.8, 70.8, 0.9 & 80.10 m$^3$ & Direct \\
\addlinespace
5 & Barrel length & $V = \pi R^2 L + \tfrac{4}{3}\pi R^2 h_{cap}$, solved for $L$ & $V$, $R$, $h_{cap}$ & 13.9675 m & Direct \\
\addlinespace
6 & Total length and slenderness & $L + 2h_{cap}$; \; $L/2R$ & $L$, $h_{cap}$, $R$ & 15.594 m; 5.37 & Direct \\
\addlinespace
7 & Shell area, exact spheroidal caps & $2\pi R L + 2\pi R^2\left[1 + \frac{1-e^2}{e}\operatorname{artanh}e\right]$ & $R$, $L$, $h_{cap}$ & 130.44 m$^2$ & Direct \\
\bottomrule
\end{tabularx}
\end{table}


\section{Stage 2: Heat Leak Budget and Insulation}
\label{sec:app_thermal}

\begin{table}[!htbp]
\centering
\caption{Heat leak budget, insulation and outer jacket, E190. Scripts: the heat leak, insulation and jacket scripts of the repository \cite{ref32}.}
\label{tab:app_thermal}
\footnotesize
\setlength{\tabcolsep}{4pt}
\begin{tabularx}{\textwidth}{@{}c>{\raggedright\arraybackslash}p{2.6cm}>{\raggedright\arraybackslash}X>{\raggedright\arraybackslash}p{2.4cm}rl@{}}
\toprule
\textbf{\#} & \textbf{Computes} & \textbf{Formula} & \textbf{Inputs} & \textbf{Result} & \textbf{Method} \\
\midrule
8 & Allowable heat leak & $Q = m_{total}(u_2 - u_1)/t_{hold}$, lever rule on the parahydrogen equation of state & $V$, $f_{max}$, 1.2 and 1.448 bar, 14.7 h & 702.6 W & Iterative \\
\addlinespace
9 & Insulation flux & $q_{MLI} = Q/A$ & 702.6 W, $A$ & 5.39 W/m$^2$ & Direct \\
\addlinespace
10 & Required total resistance & $\sum R_i = (T_{warm} - T_{cold})/Q$ & 293 K, 20 K, $Q$ & 0.3886 K/W & Direct \\
\addlinespace
11 & MLI thickness & $\ln(r_2/r_1)/(2\pi k_{eff} L_{tot}) = \sum R_i - R_{walls} - R_{film}$ & $k_{eff}$, radii & 14.93 mm & Iterative \\
\addlinespace
12 & Foam-only thickness, for comparison & same, with $k = 0.0163$ W/m$\cdot$K & \cite{ref29} & 1119 mm & Iterative \\
\addlinespace
13 & Vacuum-only emissivity needed & $q = \sigma\varepsilon_{\mathrm{eff}}(T_{warm}^4 - T_{cold}^4)$ & $q$ & 0.025 each face & Direct \\
\addlinespace
14 & Outer jacket, unstiffened and ring-stiffened & collapse pressure of Windenburg and Trilling \cite{ref33}, factor 1.5 on 1 atm & $E$, $\nu$, jacket radius, ring spacing & 25.7; 3.45 mm & Iterative \\
\bottomrule
\end{tabularx}
\end{table}


\section{Stage 3: Loads}
\label{sec:app_forces}

\begin{table}[!htbp]
\centering
\caption{Mass build-up and loads, E190, LC7 Emergency at $\phi_f = 0.90$. Script: E190\_MASTER\_thermo\_structural.py.}
\label{tab:app_forces}
\footnotesize
\setlength{\tabcolsep}{4pt}
\begin{tabularx}{\textwidth}{@{}c>{\raggedright\arraybackslash}p{2.6cm}>{\raggedright\arraybackslash}X>{\raggedright\arraybackslash}p{2.4cm}rl@{}}
\toprule
\textbf{\#} & \textbf{Computes} & \textbf{Formula} & \textbf{Inputs} & \textbf{Result} & \textbf{Method} \\
\midrule
15 & Shell mass & $m_{shell} = \rho A t$ & 2850, $A$, 1.289 mm & 479 kg & Direct \\
\addlinespace
16 & Liquid mass at $\phi_f$ & $m_{liq} = \phi_f\,\rho_{LH2}V$ & 0.9, 70.8, $V$ & 5104 kg & Direct \\
\addlinespace
17 & Resultant load factor and tilt & $\lvert\mathbf{n}\rvert = \sqrt{n_x^2+n_y^2+n_z^2}$; $\cos\theta = n_z/\lvert\mathbf{n}\rvert$ & 9, 3, 6 & 11.22; $57.7^\circ$ & Direct \\
\addlinespace
18 & Applied resultant & $F = (m_{shell}+m_{liq})\,\lvert\mathbf{n}\rvert\,g$ & masses, 11.22, 9.81 & 614.8 kN & Direct \\
\addlinespace
19 & Free surface & plane normal to $\hat{\mathbf{g}}$ enclosing $\phi_f V$, by sorted volume samples & $\hat{\mathbf{g}}$, $\phi_f$ & depth 11.45 m & Iterative \\
\addlinespace
20 & Wall pressure & $p + \rho_{LH2}\lvert\mathbf{n}\rvert g\max(0, s - s_{surf})$ & 1.5 bar, depth & up to 2.392 bar & Direct \\
\bottomrule
\end{tabularx}
\end{table}


\section{Stage 4: Structural Evaluation and Sizing}
\label{sec:app_structural}

Entry 21 is the only structural quantity not produced by a formula; every check downstream consumes it.

\begin{table}[!htbp]
\centering
\caption{Structural checks at the sized wall, E190, LC7 Emergency at $\phi_f = 0.90$. Scripts: E190\_MASTER\_thermo\_structural.py, final\_numbers.py.}
\label{tab:app_structural}
\footnotesize
\setlength{\tabcolsep}{4pt}
\begin{tabularx}{\textwidth}{@{}c>{\raggedright\arraybackslash}p{2.6cm}>{\raggedright\arraybackslash}X>{\raggedright\arraybackslash}p{2.4cm}rl@{}}
\toprule
\textbf{\#} & \textbf{Computes} & \textbf{Formula} & \textbf{Inputs} & \textbf{Result} & \textbf{Method} \\
\midrule
21 & Temperature and stress fields & discretised heat balance, then $\mathbf{K}\mathbf{u} = \mathbf{F}$ & mesh, loads, rings & fields & FEM \\
\addlinespace
22 & Membrane and bending parts & $\tfrac{1}{2}(\sigma_{top} \pm \sigma_{bot})$ & section-point stresses & per element & Direct \\
\addlinespace
23 & Allowable stress & $\sigma_{allow} = F_{ty}/\mathrm{SF}$ & 393, 1.5 & 262.0 MPa & Direct \\
\addlinespace
24 & Yield factor of safety, \SI{1.5}{\bar} & $\sigma_{allow}/\sigma_{vM}$, far field & 262.0, 256.93 MPa & 1.020 & Direct \\
\addlinespace
25 & Classical buckling stress & $\sigma_{cl} = E t/\big(R\sqrt{3(1-\nu^2)}\big)$ & 73.1 GPa, 1.289 mm, 1.3015 m & 44.27 MPa & Direct \\
\addlinespace
26 & SP-8007 knockdown \cite{ref31} & $\phi = \sqrt{R/t}/16$; $\gamma = 1 - 0.901(1 - e^{-\phi})$ & $R$, $t$ & 1.986; 0.2226 & Direct \\
\addlinespace
27 & Buckling allowable & $\sigma_{cr} = \gamma\sigma_{cl}$ & $\gamma$, $\sigma_{cl}$ & 9.85 MPa & Direct \\
\addlinespace
28 & Buckling factor of safety, \SI{0.187}{\bar} & $\sigma_{cr}/\lvert\sigma_{x,\min}\rvert$, compression only & 9.85, -5.95 MPa & 1.657 & Direct \\
\addlinespace
29 & Sized wall & search on $t$ until the governing factor lies within 6 \% above its target & entries 24 and 28 & 1.289 mm & Iterative (FEM in loop) \\
\bottomrule
\end{tabularx}
\end{table}

The same entries for the other two aircraft give $\sigma_{cr}$ = 8.80 MPa and 33.16 MPa, and the walls, governing requirements and factors of Table~\ref{tab:sized_walls}.


\section{Stage 5: Cooldown Reference Values}
\label{sec:app_thermal_ref}

\begin{table}[!htbp]
\centering
\caption{Cooldown reference values, E190. These are not inputs to the analysis, whose stress-free state is the cold state. Script: hand\_calcs.py.}
\label{tab:app_thermal_ref}
\footnotesize
\setlength{\tabcolsep}{4pt}
\begin{tabularx}{\textwidth}{@{}c>{\raggedright\arraybackslash}p{2.6cm}>{\raggedright\arraybackslash}X>{\raggedright\arraybackslash}p{2.4cm}rl@{}}
\toprule
\textbf{\#} & \textbf{Computes} & \textbf{Formula} & \textbf{Inputs} & \textbf{Result} & \textbf{Method} \\
\midrule
30 & Cooldown temperature drop & $\Delta T = T_{asm} - T_{liq}$ & 293, 20.3 & 272.7 K & Direct \\
\addlinespace
31 & Fully restrained stress, uniaxial and biaxial & $E\varepsilon_{c}$; \; $E\varepsilon_{c}/(1-\nu)$ & $E$, $\varepsilon_{c} = 0.415\%$ \cite{ref34}, $\nu$ & 304; 453 MPa & Direct \\
\addlinespace
32 & Free axial contraction & $\varepsilon_{c}\,L_{tot}$ & $\varepsilon_{c}$, 15.594 m & 64.8 mm & Direct \\
\bottomrule
\end{tabularx}
\end{table}

Entries 31 and 32 bracket the cooldown problem, using the room-temperature expansion coefficient held constant over the whole drop. Because the coefficient of aluminium falls steeply at low temperature, both values overstate the true ones; they are quoted only to show why the cold state is taken as stress free and why the supports must let the vessel contract towards its anchor.


\section{Verification Against Hand Calculation}
\label{sec:app_verification}

Table~\ref{tab:app_agreement} sets every closed-form check of Section~\ref{sec:verification} beside the finite element value, for all three aircraft at their sized walls. The closed-form values use no finite element input, so each agreement is independent evidence. Script: hand\_calcs.py \cite{ref32}.

\begin{table}[!htbp]
\centering
\caption[Hand calculation against the finite element model]{Hand calculation against the finite element model at the sized walls, given as closed form / model. Baseline, $\phi_f = 0.10$, unless stated; the waterline flux is at $\phi_f = 0.90$.}
\label{tab:app_agreement}
\footnotesize
\setlength{\tabcolsep}{4pt}
\begin{tabular}{@{}lccc@{}}
\toprule
\textbf{Check} & \textbf{E190} & \textbf{A320} & \textbf{A350-900} \\
\midrule
Barrel hoop stress, $pR/t$ [MPa] & 151.50 / 151.46 & 163.01 / 162.95 & 68.13 / 68.05 \\
Mean barrel axial stress, $pR/2t$ [MPa] & 75.75 / 75.75 & 81.51 / 81.51 & 34.06 / 34.07 \\
Cap equator hoop stress [MPa] & -42.4 / -51.0 & -45.6 / -53.2 & -19.1 / -25.6 \\
Equilibrium, LC7, $\phi_f = 0.90$ [kN] & 614.8 / 614.2 & 802.1 / 801.1 & 6354.8 / 6348.5 \\
Dry-wall $\Delta T = q/h_{dry}$ [K] & 3.689 / 3.683 & 4.651 / 4.644 & 7.376 / 7.364 \\
Waterline flux $q L_{fin}/t$ [W/m$^2$] & 1338 / 1317 & 1527 / 1489 & 1274 / 1284 \\
Mean axial rise, LC7 minus Baseline, $\phi_f = 0.90$ [MPa] & 21.38 / 21.04 & 17.48 / 17.08 & 22.98 / 22.67 \\
Min.\ axial stress, LC7, \SI{0.187}{\bar} [MPa] & -11.52 / -5.95 & 1.57 / 1.71 & -26.72 / -20.85 \\
Mid-ring reaction from bowing [kN] & -2.5 / -1.8 & -24.5 / -11.4 & -33.3 / -29.8 \\
Bending decay length $\pi/\beta$ [m] & 0.101 / - & 0.127 / - & 0.297 / - \\
Fin length $L_{fin}$ [m] & 0.320 / - & 0.353 / - & 0.671 / - \\
\bottomrule
\end{tabular}
\end{table}


\section{Complete Results Matrices}
\label{sec:app_full_matrices}

The complete production and stability matrices of all three aircraft are listed here, so that every one of the 120 cases appears at row level. Columns: wetted-to-dry temperature difference, far-field von Mises stress and yield factor of safety at \SI{1.5}{\bar}, maximum displacement and total ring reaction at \SI{1.5}{\bar}, and the most compressive axial membrane stress in the barrel and the buckling factor of safety at \SI{0.187}{\bar} (a dash where the barrel is in tension). The governing case is in bold. Source: production\_matrix\_sized\_walls.csv and final\_numbers.json \cite{ref32}.


\begin{table}[!htbp]
\centering
\caption[Complete results matrix for the E190]{Complete results matrix for the E190, wall 1.289 mm, 238,169 elements.}
\label{tab:app_e190_full}
\footnotesize
\setlength{\tabcolsep}{4pt}
\begin{tabular}{@{}llrrrrrrr@{}}
\toprule
 & & \multicolumn{5}{c}{\textbf{\SI{1.5}{\bar}}} & \multicolumn{2}{c}{\textbf{\SI{0.187}{\bar}}} \\
\cmidrule(lr){3-7}\cmidrule(l){8-9}
\textbf{Case} & $\phi_f$ & $\Delta T$ [K] & $\sigma_{vM}$ [MPa] & FOS & $u_{\max}$ [mm] & $R$ [kN] & $\sigma_{x,\min}$ [MPa] & FOS$_b$ \\
\midrule

Baseline & 0.10 & 3.683 & 162.6 & 1.611 & 7.77 & 10.1 & 6.90 & - \\
Baseline & 0.30 & 3.681 & 162.8 & 1.610 & 7.68 & 21.3 & 6.08 & - \\
Baseline & 0.50 & 3.672 & 163.2 & 1.606 & 7.53 & 32.4 & 5.47 & - \\
Baseline & 0.75 & 3.619 & 163.3 & 1.605 & 7.48 & 46.4 & 4.52 & - \\
Baseline & 0.90 & 3.432 & 163.5 & 1.602 & 7.45 & 54.7 & 4.07 & - \\
\addlinespace
LC1 & 0.10 & 3.683 & 162.7 & 1.610 & 7.77 & 25.4 & 6.02 & - \\
LC1 & 0.30 & 3.681 & 162.9 & 1.608 & 7.68 & 53.2 & 3.78 & - \\
LC1 & 0.50 & 3.672 & 163.3 & 1.604 & 7.52 & 81.0 & 2.46 & - \\
LC1 & 0.75 & 3.619 & 164.0 & 1.597 & 7.50 & 115.9 & 2.01 & - \\
LC1 & 0.90 & 3.432 & 164.6 & 1.592 & 7.49 & 136.8 & 2.34 & - \\
\addlinespace
LC3 & 0.10 & 3.684 & 164.8 & 1.590 & 7.92 & 29.9 & 8.63 & - \\
LC3 & 0.30 & 3.684 & 168.4 & 1.556 & 7.92 & 62.3 & 7.36 & - \\
LC3 & 0.50 & 3.684 & 171.5 & 1.528 & 7.91 & 94.6 & 6.77 & - \\
LC3 & 0.75 & 3.684 & 175.8 & 1.491 & 8.06 & 135.3 & 5.77 & - \\
LC3 & 0.90 & 3.675 & 178.4 & 1.469 & 8.18 & 159.6 & 4.47 & - \\
\addlinespace
LC7 & 0.10 & 3.684 & 175.5 & 1.492 & 8.77 & 115.3 & 6.84 & - \\
LC7 & 0.30 & 3.684 & 195.9 & 1.338 & 9.83 & 239.5 & 6.83 & - \\
LC7 & 0.50 & 3.684 & 216.1 & 1.212 & 10.68 & 364.5 & 6.68 & - \\
LC7 & 0.75 & 3.684 & 241.7 & 1.084 & 11.82 & 520.7 & 0.61 & - \\
\textbf{LC7} & \textbf{0.90} & \textbf{3.678} & \textbf{256.9} & \textbf{1.020} & \textbf{12.50} & \textbf{614.2} & \textbf{-5.95} & \textbf{1.66} \\
\bottomrule
\end{tabular}
\end{table}


\begin{table}[!htbp]
\centering
\caption[Complete results matrix for the A320]{Complete results matrix for the A320, wall 1.572 mm, 165,602 elements.}
\label{tab:app_a320_full}
\footnotesize
\setlength{\tabcolsep}{4pt}
\begin{tabular}{@{}llrrrrrrr@{}}
\toprule
 & & \multicolumn{5}{c}{\textbf{\SI{1.5}{\bar}}} & \multicolumn{2}{c}{\textbf{\SI{0.187}{\bar}}} \\
\cmidrule(lr){3-7}\cmidrule(l){8-9}
\textbf{Case} & $\phi_f$ & $\Delta T$ [K] & $\sigma_{vM}$ [MPa] & FOS & $u_{\max}$ [mm] & $R$ [kN] & $\sigma_{x,\min}$ [MPa] & FOS$_b$ \\
\midrule

Baseline & 0.10 & 4.644 & 178.9 & 1.464 & 9.21 & 13.1 & 7.32 & - \\
Baseline & 0.30 & 4.643 & 179.0 & 1.464 & 9.12 & 27.7 & 6.80 & - \\
Baseline & 0.50 & 4.640 & 179.4 & 1.461 & 8.89 & 42.2 & 6.44 & - \\
Baseline & 0.75 & 4.611 & 179.7 & 1.458 & 8.86 & 60.4 & 5.76 & - \\
Baseline & 0.90 & 4.467 & 180.1 & 1.455 & 8.85 & 71.4 & 4.77 & - \\
\addlinespace
LC1 & 0.10 & 4.644 & 178.9 & 1.464 & 9.21 & 32.7 & 7.06 & - \\
LC1 & 0.30 & 4.643 & 179.0 & 1.464 & 9.11 & 69.2 & 6.40 & - \\
LC1 & 0.50 & 4.640 & 179.3 & 1.461 & 8.88 & 105.5 & 5.99 & - \\
LC1 & 0.75 & 4.611 & 180.6 & 1.451 & 8.90 & 151.1 & 5.33 & - \\
LC1 & 0.90 & 4.467 & 181.6 & 1.443 & 8.92 & 178.5 & 4.42 & - \\
\addlinespace
LC3 & 0.10 & 4.644 & 179.3 & 1.461 & 9.32 & 38.5 & 9.48 & - \\
LC3 & 0.30 & 4.644 & 182.8 & 1.433 & 9.30 & 81.0 & 9.03 & - \\
LC3 & 0.50 & 4.644 & 185.5 & 1.412 & 9.26 & 123.2 & 7.60 & - \\
LC3 & 0.75 & 4.644 & 189.0 & 1.386 & 9.37 & 176.3 & 5.46 & - \\
LC3 & 0.90 & 4.630 & 191.5 & 1.368 & 9.49 & 208.3 & 6.93 & - \\
\addlinespace
LC7 & 0.10 & 4.644 & 189.2 & 1.385 & 10.01 & 148.1 & 8.63 & - \\
LC7 & 0.30 & 4.644 & 206.3 & 1.270 & 10.88 & 311.1 & 8.63 & - \\
LC7 & 0.50 & 4.644 & 223.2 & 1.174 & 11.63 & 474.1 & 7.42 & - \\
LC7 & 0.75 & 4.644 & 244.6 & 1.071 & 12.65 & 679.5 & 3.89 & - \\
\textbf{LC7} & \textbf{0.90} & \textbf{4.631} & \textbf{257.3} & \textbf{1.018} & \textbf{13.27} & \textbf{801.1} & \textbf{1.71} & \textbf{-} \\
\bottomrule
\end{tabular}
\end{table}


\begin{table}[!htbp]
\centering
\caption[Complete results matrix for the A350-900]{Complete results matrix for the A350-900, wall 5.674 mm, 126,816 elements.}
\label{tab:app_a350_full}
\footnotesize
\setlength{\tabcolsep}{4pt}
\begin{tabular}{@{}llrrrrrrr@{}}
\toprule
 & & \multicolumn{5}{c}{\textbf{\SI{1.5}{\bar}}} & \multicolumn{2}{c}{\textbf{\SI{0.187}{\bar}}} \\
\cmidrule(lr){3-7}\cmidrule(l){8-9}
\textbf{Case} & $\phi_f$ & $\Delta T$ [K] & $\sigma_{vM}$ [MPa] & FOS & $u_{\max}$ [mm] & $R$ [kN] & $\sigma_{x,\min}$ [MPa] & FOS$_b$ \\
\midrule

Baseline & 0.10 & 7.364 & 61.3 & 4.273 & 9.08 & 148.3 & -0.94 & 35.25 \\
Baseline & 0.30 & 7.356 & 62.0 & 4.228 & 8.63 & 252.7 & -2.78 & 11.93 \\
Baseline & 0.50 & 7.330 & 62.2 & 4.213 & 8.06 & 356.9 & -4.20 & 7.90 \\
Baseline & 0.75 & 7.199 & 61.9 & 4.235 & 7.68 & 488.0 & -5.62 & 5.90 \\
Baseline & 0.90 & 6.771 & 61.9 & 4.230 & 7.46 & 566.3 & -6.81 & 4.87 \\
\addlinespace
LC1 & 0.10 & 7.364 & 64.1 & 4.085 & 9.07 & 370.7 & -1.37 & 24.28 \\
LC1 & 0.30 & 7.356 & 65.9 & 3.974 & 9.04 & 631.8 & -3.15 & 10.52 \\
LC1 & 0.50 & 7.330 & 65.7 & 3.988 & 9.44 & 892.3 & -4.51 & 7.35 \\
LC1 & 0.75 & 7.199 & 66.0 & 3.967 & 9.30 & 1219.8 & -5.52 & 6.01 \\
LC1 & 0.90 & 6.771 & 64.0 & 4.093 & 9.32 & 1415.8 & -7.37 & 4.50 \\
\addlinespace
LC3 & 0.10 & 7.366 & 63.1 & 4.151 & 9.80 & 434.6 & 2.23 & - \\
LC3 & 0.30 & 7.366 & 66.2 & 3.956 & 9.80 & 737.6 & 0.62 & - \\
LC3 & 0.50 & 7.366 & 69.3 & 3.781 & 9.49 & 1041.6 & -0.95 & 34.72 \\
LC3 & 0.75 & 7.366 & 73.1 & 3.585 & 8.71 & 1421.6 & -1.57 & 21.16 \\
LC3 & 0.90 & 7.350 & 75.4 & 3.475 & 8.97 & 1649.6 & -3.67 & 9.03 \\
\addlinespace
LC7 & 0.10 & 7.366 & 72.5 & 3.613 & 11.71 & 1673.9 & -3.71 & 8.95 \\
LC7 & 0.30 & 7.366 & 90.8 & 2.886 & 13.62 & 2837.9 & -3.71 & 8.93 \\
LC7 & 0.50 & 7.366 & 109.0 & 2.404 & 15.06 & 4009.6 & -2.32 & 14.32 \\
LC7 & 0.75 & 7.366 & 132.0 & 1.985 & 17.58 & 5473.4 & -15.22 & 2.18 \\
\textbf{LC7} & \textbf{0.90} & \textbf{7.358} & \textbf{145.7} & \textbf{1.798} & \textbf{19.12} & \textbf{6348.5} & \textbf{-20.85} & \textbf{1.59} \\
\bottomrule
\end{tabular}
\end{table}
